\section{Evaluation Methodology}
\label{sec:evaluation_setup}

The evaluation addresses a particular sub-problem within the field of automatic assessment in education~\cite{ref_ihantola2010,ref_paiva2022}, namely semantic strategy retrieval for multiple programming languages in the CS-1 context. All similarity evaluation design decisions are compatible with those made in Phase~1 in order to allow for cross-phase performance comparison. Generative evaluation is another independent evaluation effort focused on the same 20 CS-1 algorithmic problems. All scope caveats are summarized in Section~\ref{subsec:threats}; the results are discussed in terms of CS-1 scope in Section~\ref{subsec:scope}.

\subsection{Dataset Construction}
\label{subsec:dataset_construction}

\textbf{Problem set:} The evaluation involves 20 CS-1 imperative algorithmic problems: linear and binary search, bubble sort, selection sort, insertion sort, sum of array, minimum and maximum element, GCD/LCM, factorial, Fibonacci (iterative and recursive), palindrome check, prime check, string reversal, digit sum, power computation, matrix addition, array rotation, and common element detection. These problems collectively exercise all 21 CES taxonomy patterns and span the full CS-1 imperative algorithmic scope.

% \textbf{Monolingual cohorts:} Two monolingual cohorts were assembled from student C++ and Java submissions at Amrita Vishwa Vidyapeetham, Amritapuri, India (20 problems $\times$ 20 submissions per language), collected during a different academic semester than the cross-language corpus. Reference solutions for each problem were authored by an instructor independent of the research team.

\textbf{Cross-language corpus and filtering criteria:} The 400-program cross-language corpus consists of student submissions written in C, C++, and Java collected from an introductory CS-1 course. Submissions were filtered based on compilability and valid entry-point identification using a curated 23-pattern entry-point list, ensuring structural integrity for Joern CPG graph generation and pattern extraction. Ground truth strategy labels for all 400 programs were assigned under the strict hard-filtering constraint (Section~\ref{subsec:late_fusion}).

\textbf{APM generative corpus:} Program pairs make up a full 6-directional matrix (C$\rightarrow$C++, C$\rightarrow$Java, C++$\rightarrow$C, C++$\rightarrow$Java, Java$\rightarrow$C, Java$\rightarrow$C++) with 120 program pairs generated across 20 different problems for generative APM evaluation.

\textbf{Total corpus scale:} Across all evaluation tracks, the study involves 1,320 distinct program instances: 400 C++ (monolingual), 400 Java (monolingual), 400 cross-language (C/C++/Java), and 120 APM generation targets. Given that the CS-1 domain assumes a finite set of canonical strategies for any algorithmic problem (Section~\ref{subsec:ces_stability}), the 400-program data set offers enough statistical precision for 95\% Wilson Confidence Intervals. This is confirmed by 5-fold cross-validation of semantic weights on C++ and Java monolingual corpora (Section~\ref{sec:results}).

\subsection{Ground Truth Annotation Protocol}
\label{subsec:ground_truth}

Each individual student submission within the cross-language data collection received an associated ground-truth \textit{dominant algorithmic strategy} classification taken from a closed set of 20 canonical algorithms (one per problem). Top-1 retrievals are considered correct if the most similar retrieved reference solution bears the same classification for dominant algorithmic strategy as the query submission and is from a language distinct from that of the query (hard-filtered constraint).


\textbf{Annotation:} Annotations were performed by the lead author for all $N=400$ submissions. To validate the reliability of this procedure, an Inter-Rater Reliability (IRR) study was carried out for a stratified random sample of $n=60$ submissions independently by a second domain expert not involved in the taxonomy development. This produced an inter-annotator agreement of Cohen's $\kappa = 0.86$~\cite{ref_cohen1960} ($p < 0.001$), which suggests a strong level of agreement according to Landis and Koch~\cite{ref_landis1977}. Conflict of interests of the lead author are noted and addressed in Threat Group~1 of Section~\ref{subsec:threats}.

\textbf{APM evaluation:} Compilation and behavioral agreement are evaluated for each of the $N=120$ generated APM programs. Behaviorally, each program should agree with the specification given the test cases consisting of three to five inputs within the respective direction. Such a test suite was created by the lead author. This design constraint and its possible effects are considered in Threat Group~5 of Section~\ref{subsec:threats}.

\subsection{UniXcoder Fine-Tuning Protocol}
\label{subsec:neural_finetuning}

For the purpose of providing a neural comparison point besides the zero-shot results, the UniXcoder~\cite{ref_unixcoder} model was fine-tuned on a subset of our dataset using the following protocol:

\textbf{Train-test split:} The cross-language dataset containing 400 programs was split by stratified sampling into 80/20 train-test ratio, leading to 320 programs used for fine-tuning and 80 programs for evaluation. The evaluation set of 80 programs is completely disjoint with respect to the fine-tuning data: there is no intersection between them.

\textbf{Fine-tuning procedure:} UniXcoder was fine-tuned based on the InfoNCE loss function on pairs of programs with the same solution strategy within a cross-language setup selected from 320 training programs. The fine-tuning process was executed with AdamW optimizer, a learning rate equal to $2 \times 10^{-5}$, batch size equals to 16, the number of training epochs being 10 and a random seed of 42. Fine-tuning was performed on an NVIDIA A100 GPU using PyTorch version 2.1.0 (CUDA 11.8) with Hugging Face \texttt{transformers} package version 4.35.0 and took around 14 minutes. The reported accuracy of 71.37\% represents the performance of a single training run initialized with a fixed random seed (seed 42) to ensure reproducibility, serving as a baseline comparison.

\textbf{Evaluation:} Evaluation of the fine-tuned UniXcoder was done on the test set of 80 programs with the help of cosine similarity between CLS tokens in top-1 retrieval mode according to the hard filter criterion. The result of 71.37\% [95\% Wilson CI: 62.15\%, 79.12\%] corresponds to 57/80 correct retrievals. The wider CI relative to the full-dataset CIs reflects the smaller evaluation set size ($n=80$ vs. $n=400$) and is expected.

\textbf{Fair comparison:} The aggregate performance of the framework (87.25\%) is measured using the whole set of 400 programs. To ensure an unbiased comparison, the proposed framework was tested on a held-out set of 80 programs yielding 88.75\% [95\% Wilson CI: 80.49\%, 93.83\%] (71 out of 80 correct), which demonstrates the result of 17.38 percentage points improvement on the common set of data. While conducting a fair comparison, it should be noted that the fine-tuning protocol relies on training on 320 programs, which is significantly fewer examples in comparison to a practical setting. Meanwhile, a deterministic framework does not rely on any labeled data.

\subsection{CodeBERT Fine-Tuning Protocol}
\label{subsec:codebert_finetuning}

To provide another example of fine-tuned neural network models and prove that the observed improvement for fine-tuning UniXcoder is not model-specific, the CodeBERT model~\cite{ref_codebert} was fine-tuned independently on the same dataset using the same protocol:

\textbf{Train/test split:} Identical 80/20 stratified split used in the UniXcoder fine-tuning protocol (Section~\ref{subsec:neural_finetuning}), 320 fine-tuning programs and 80-program held-out test set. The held-out set for evaluation is the same for two fine-tuned neural models; the fine-tuning sets do not intersect.

\textbf{Fine-tuning procedure:} CodeBERT was fine-tuned using the same InfoNCE loss (with temperature $\tau=0.07$) objective but on same-strategy cross-language program pairs from the 320-programs training set. Training used the AdamW optimizer, learning rate $2\times10^{-5}$, batch size of 8, 3 training epochs and seed 42. Max tokenized program length was set to 128 tokens (instead of 512, due to CPU constraints). Fine-tuning was performed on CPU with PyTorch and the Hugging Face \texttt{transformers}. The choice of 3-epochs fine-tune is an intentional CPU-compatible version of the 10-epochs GPU fine-tune procedure used for UniXcoder. This distinction is important, since the difference in fine-tuning time is indicated in Table~\ref{tab:baseline_comparison}. 

\textbf{Evaluation:} The fine-tuned model was evaluated using the same set of 80 held-out programs, using CLS-token based cosine similarity with hard-filtering criterion for top-1 retrieval. The results are shown in Table~\ref{tab:baseline_comparison} ($^{\S}$ footnote).



\subsection{Evaluation Metrics}
\label{subsec:metrics}

\textbf{Top-1 Retrieval Accuracy:} The main metric is the top-1 retrieval accuracy---the ratio of query submissions in which the top-ranked reference program belongs to the same strategy category. All accuracy values are reported in $n/N$ counts and 95\% Wilson confidence intervals.

\textbf{McNemar Test:} Paired comparison between the proposed framework and the TF-IDF baseline via the McNemar test, which operates on per-submission binary outcomes. The Yates-corrected McNemar test statistic is defined as:
\begin{equation}
\chi^2 = \frac{(|b - c| - 0.5)^2}{b + c}
\end{equation}
where $b$ is the number of queries correct only in the proposed framework, and $c$ is the number of queries correct only in the baseline. A contingency table 2$\times$2 is provided (see Table~\ref{tab:mcnemar_contingency} for exact numbers); independent verification can be done via this test since the same 400 submissions were evaluated in both conditions (the proposed framework vs. TF-IDF baseline).

\textbf{Two-Proportion Z-Test:} Comparison between C++ monolingual and cross-language cohorts through the two-proportion z-test (also see Section~\ref{subsec:weight_stability_discussion}). The test statistic is given by:
\begin{equation}
z = \frac{\hat{p}_1 - \hat{p}_2}{\sqrt{\hat{p}(1-\hat{p})\left(\frac{1}{n_1} + \frac{1}{n_2}\right)}}
\end{equation}
where $\hat{p}_1$ and $\hat{p}_2$ are the sample proportions of correct retrievals, $n_1$ and $n_2$ are the group sizes, and $\hat{p}$ is the pooled success proportion:
\begin{equation}
\hat{p} = \frac{n_1\hat{p}_1 + n_2\hat{p}_2}{n_1 + n_2}
\end{equation}
This test is valid since there are two different program instances involved in the comparison: the C++ monolingual cohort of $n=400$ student C++ programs and the cross-language cohort of $n=400$ programs ($n_{c}=130$ C, $n_{cpp}=130$ C++, and $n_{java}=140$ Java) that are obtained from different student cohorts within the 20 problem institution corpus. Population independence was confirmed by verifying submission IDs; no instance of program appeared in both evaluation sets.

\textbf{Odds Ratio:} OR $= (a \times d)/(b \times c)$ from the McNemar contingency table; 95\% CI computed as $\exp(\ln(\text{OR}) \pm 1.96 \cdot \sqrt{1/a + 1/b + 1/c + 1/d})$.

\textbf{Wilson Confidence Intervals:} All proportions reported as $\hat{p}$ [95\% Wilson CI: $p_L$, $p_U$] where:
\begin{equation}
p_{L,U} = \frac{\hat{p} + \frac{z^2}{2n} \pm z\sqrt{\frac{\hat{p}(1-\hat{p})}{n} + \frac{z^2}{4n^2}}}{1 + \frac{z^2}{n}}, \quad z = 1.96
\end{equation}

\textbf{5-Fold Cross-Validation:} Conducted on both monolingual cohorts to check whether weights generalize to unseen languages within language. Cross-validation did not take place for cross-language cohort; two-proportion z-test (Section~\ref{subsec:weight_stability_discussion}) and weight sensitivity sweep (see Table~\ref{tab:tversky_sensitivity}) are primary evidence of cross-language weight stability. Performing cross-validation for cross-language cohort remains future work, described as validity threat in Section~\ref{subsec:threats}.

\textbf{Statistical Test Assumptions:} \textit{McNemar Test:} (i)~The paired binary outcomes are measured on the same sample of $N{=}400$ student submissions evaluated under two distinct conditions (proposed framework vs. TF-IDF baseline); (ii)~the test statistic relies on the counts of discordant pairs (cells $b$ and $c$ in the contingency table); and (iii)~Yates-corrected $\chi^2$ values are reported alongside raw values to ensure a conservative assessment. \textit{Two-Proportion Z-Test:} (i)~The two groups being compared are independent (confirmed by verifying that submission IDs are completely disjoint between the monolingual cohort and the cross-language cohort); (ii)~the outcome of each retrieval is a binary event (correct/incorrect retrieval); and (iii)~the sample size is sufficiently large to satisfy the normal approximation condition ($n\hat{p} > 5$ and $n(1-\hat{p}) > 5$ for all groups under comparison).

